% ===================================================================
%  Generado por Erlen
%  Compílalo en Overleaf, o localmente con dos pasadas de pdflatex.
% ===================================================================
\documentclass[aspectratio=169,11pt]{beamer}
\usetheme{Boadilla}
\usepackage{lmodern}
% La letra de Beamer ya es Latin Modern Sans, la misma que ves en pantalla: aquí solo se tiñe el papel.
\definecolor{arenaPapel}{HTML}{F2EADB}
\definecolor{arenaTinta}{HTML}{33291F}
\definecolor{arenaCobre}{HTML}{8A571C}
\setbeamercolor{background canvas}{bg=arenaPapel}
\setbeamercolor{normal text}{fg=arenaTinta}
\setbeamercolor{structure}{fg=arenaCobre}
% El doble filete bajo el título, uno grueso de cobre y otro fino, es cosa de Erlen: en Beamer habría que redefinir la plantilla del frametitle.
\usepackage[utf8]{inputenc}
\usepackage[T1]{fontenc}
% Español: se carga solo si babel-spanish está instalado (Overleaf lo trae).
\IfFileExists{spanish.ldf}{\usepackage[spanish,es-nodecimaldot]{babel}}{}
\usepackage{amsmath,amssymb}
\usepackage[version=4]{mhchem}
\usepackage{booktabs}
\definecolor{erlenfondo}{HTML}{F2EADB}
\definecolor{erlenacento}{HTML}{8A571C}
\usepackage{tikz}
\usetikzlibrary{arrows.meta, calc}
\renewcommand{\figurename}{Figura}
\renewcommand{\tablename}{Tabla}
\usepackage{siunitx}
\sisetup{per-mode=symbol, range-units=single}
\setbeamertemplate{navigation symbols}{}
% Algunos temas (Metropolis entre ellos) insertan solos una diapositiva al
% empezar cada sección. Erlen ya emite la suya, así que se desactiva la
% automática: de otro modo cada sección saldría dos veces y la numeración
% del PDF dejaría de coincidir con la de la app.
\AtBeginSection{}
\AtBeginSubsection{}

\title[Equilibrio ácido–base]{Equilibrio ácido–base}
\subtitle{Ejemplo didáctico · datos ilustrativos}
\date{\today}

\begin{document}

% --- diapositiva 1 ---
\begin{frame}[plain]
  \transdissolve[duration=0.25]
  \titlepage
\end{frame}
\note{%
  {\scriptsize\textbf{Tiempo previsto: 0:30 min}}\par\medskip
  Ejemplo didáctico. Sustituye contenido y datos por los de tu trabajo. No atribuir estas cifras a un experimento.\par
}

% --- diapositiva 2 ---
\begin{frame}{La pregunta que guía el trabajo}
  \transdissolve[duration=0.25]
  {\large ¿Qué predice un modelo ideal para la relación ácido/base?}
  \medskip

  {\small Caso de uso ilustrativo. Define aquí el sistema, la comparación y el alcance.}
\end{frame}
\note{%
  {\scriptsize\textbf{Tiempo previsto: 1:30 min}}\par\medskip
}

% --- diapositiva 3 ---
\begin{frame}{Una estrategia explícita}
  \transdissolve[duration=0.25]
  \begin{figure}
    \centering
    \definecolor{sa0f}{HTML}{CFDCE6}
    \definecolor{sa1f}{HTML}{9DB7CC}
    \definecolor{sa2f}{HTML}{6B93B2}
    \definecolor{sa3f}{HTML}{396E98}
    \definecolor{sat1843819360}{HTML}{15181C}
    \definecolor{sat1226267613}{HTML}{FFFFFF}
    % Con babel-spanish las palabras largas se parten solas y el texto respira mejor.
    \begin{tikzpicture}[x=1cm,y=1cm,font=\footnotesize]
      \draw[fill=sa0f, draw=none] (0.000,0.000) -- (2.226,0.000) -- (2.518,-1.254) -- (2.226,-2.508) -- (0.000,-2.508) -- cycle;
      \draw[fill=sa1f, draw=none] (2.644,0.000) -- (4.870,0.000) -- (5.162,-1.254) -- (4.870,-2.508) -- (2.644,-2.508) -- (2.936,-1.254) -- cycle;
      \draw[fill=sa2f, draw=none] (5.288,0.000) -- (7.514,0.000) -- (7.806,-1.254) -- (7.514,-2.508) -- (5.288,-2.508) -- (5.580,-1.254) -- cycle;
      \draw[fill=sa3f, draw=none] (7.932,0.000) -- (10.157,0.000) -- (10.450,-1.254) -- (10.157,-2.508) -- (7.932,-2.508) -- (8.224,-1.254) -- cycle;
      \node[text=sat1843819360, align=center, text width=1.75cm, inner sep=1pt] at (1.113,-1.254) {\bfseries Definir especies};
      \node[text=sat1843819360, align=center, text width=1.75cm, inner sep=1pt] at (3.976,-1.254) {\bfseries Plantear el equilibrio};
      \node[text=sat1843819360, align=center, text width=1.75cm, inner sep=1pt] at (6.620,-1.254) {\scriptsize \bfseries Declarar aproximaciones};
      \node[text=sat1226267613, align=center, text width=1.75cm, inner sep=1pt] at (9.264,-1.254) {\bfseries Comparar con mediciones};
    \end{tikzpicture}
  \end{figure}
\end{frame}
\note{%
  {\scriptsize\textbf{Tiempo previsto: 1:30 min}}\par\medskip
  Explica las decisiones y los controles, no solo los nombres de las técnicas.\par
}

% --- diapositiva 4 ---
\begin{frame}{Mostrar la evidencia}
  \transdissolve[duration=0.25]
  \begin{center}
    \ce{HA <=> H+ + A-}
  \end{center}
  \medskip

  \begin{equation*}
    \mathrm{pH}=\mathrm{p}K_a+\log_{10}\frac{[\mathrm{A}^{-}]}{[\mathrm{HA}]}
  \end{equation*}
  \medskip

  \begin{table}
    \centering
    \caption{Relación ideal calculada; no recomendación de preparación.}
    \begin{tabular}{rr}
      \toprule
      \textbf{[A⁻]/[HA]} & \textbf{pH − pKa} \\
      \midrule
      0.1 & −1 \\
      1 & 0 \\
      10 & 1 \\
      \bottomrule
    \end{tabular}
  \end{table}
\end{frame}
\note{%
  {\scriptsize\textbf{Tiempo previsto: 1:30 min}}\par\medskip
  Datos o modelos didácticos: reemplazar antes de presentar resultados propios.\par
}

% --- diapositiva 5 ---
\begin{frame}{Interpretación y límites}
  \transdissolve[duration=0.25]
  \begin{itemize}
    \item Las concentraciones aproximan actividades bajo supuestos explícitos.
    \item No usar esta diapositiva como protocolo de seguridad.
    \item Contrastar las aproximaciones con el sistema experimental.
  \end{itemize}
\end{frame}
\note{%
  {\scriptsize\textbf{Tiempo previsto: 1:30 min}}\par\medskip
}

% --- diapositiva 6 ---
\begin{frame}{Antes de compartir}
  \transdissolve[duration=0.25]
  \begin{itemize}
    \item Identidad de especies, temperatura y actividades
    \item Rango de validez de la aproximación
    \item Mediciones y calibración del electrodo
  \end{itemize}
  \medskip

  {\small Conserva el proyecto JSON y revisa la exportación final.}
\end{frame}
\note{%
  {\scriptsize\textbf{Tiempo previsto: 1:30 min}}\par\medskip
}

\end{document}